\documentclass{article}

% Language setting
% Replace `english' with e.g. `spanish' to change the document language
\usepackage[english]{babel}

% Set page size and margins
% Replace `letterpaper' with `a4paper' for UK/EU standard size
\usepackage[letterpaper,top=2cm,bottom=2cm,left=3cm,right=3cm,marginparwidth=1.75cm]{geometry}

% Useful packages
\usepackage{amsmath}
\usepackage{amssymb}
\usepackage{amsthm}
\usepackage{tikz}
\usepackage{graphicx}
\usepackage{xparse}
\usepackage{chngcntr}
\usepackage{mathtools}
\counterwithin*{subsection}{section}
\renewcommand{\thesubsection}{\thesection.\alph{subsection}}
\usepackage[colorlinks=true, allcolors=blue]{hyperref}
\usepackage[utf8]{inputenc}



\usepackage{cite}
\usepackage{amsmath,amssymb,amsfonts}
\usepackage{algorithmic}
\usepackage{graphicx}
\usepackage{textcomp}
\usepackage{xcolor}
\def\BibTeX{{\rm B\kern-.05em{\sc i\kern-.025em b}\kern-.08em
    T\kern-.1667em\lower.7ex\hbox{E}\kern-.125emX}}

\usepackage{hyperref}
\usepackage{url}
\usepackage{icomma}
\usepackage{soul}
\usepackage{subcaption}
\usepackage{graphicx}
\usepackage{multirow}
\usepackage{multicol}
\usepackage{caption}
\usepackage{algorithm}
\usepackage{algorithmic}
\renewcommand{\theHalgorithm}{\arabic{algorithm}}
% For theorems and such
\usepackage{amsmath}
\usepackage{amssymb}
\usepackage{mathtools}
\usepackage{amsthm}
\usepackage{wrapfig}
\usepackage{booktabs} 
\usepackage{enumitem}

\usepackage{pdfpages}

% \usepackage[bottom]{footmisc}
\usepackage{dblfloatfix}
\usepackage{indentfirst}
\usepackage{fancyhdr}
\pagestyle{fancy}


\fancyhead[R]{Steven Golob}

% \title{scratch document}
% \author{Steven Golob}

\begin{document}
% \maketitle
What is the theoretical limit to the success of a membership inference attack on DP-synthetic data?


\vspace{20pt}
\textbf{Differential Privacy}: 

Let $\mathcal{A}$ be a synthetic data generation (SDG) algorithm, then  $\mathcal{A}$ is $\epsilon$-DP iff, for all neighboring datasets $D$ and $D'$ and any subset $\mathcal{O}$ of the range of $\mathcal{A}$:

\begin{equation}\label{DPEQ}
    P[\mathcal{A}(D) \in \mathcal{O}] \leq e^\epsilon \cdot P[\mathcal{A}(D') \in \mathcal{O}]
\end{equation}

$D$ and $D'$ are real datasets that $\mathcal{A}$ is applied to, to generate synthetic datasets $\mathcal{A}(D)$ and $\mathcal{A}(D')$. The fact that inequality (\ref{DPEQ}) holds for all neighboring datasets $D$ and $D'$ and for any subset $\mathcal{O}$ of the range of $\mathcal{A}$, implies that (\ref{DPEQ}) holds in particular for 
\begin{itemize}
\item say how you pick $D$ and $D'$; differ in one record $t$
\item say how you pick $\mathcal{O}$
\end{itemize}

Let's say that some target $t \notin D$ and $t \in D'$. What can we consider $\mathcal{O}$ to be?

Let's say that $\mathcal{O}$ is the set of all datasets that, when given as input to some arbitrary attacker, would cause that attacker to classify $t$ as a member of the training data (i.e. the attacker decides that $t \in D_{train}$). Conversely, any data set not in $\mathcal{O}$ would cause the attacker to give a ``no'', when given to the attacker as input.

If this is an adequate way of thinking about $\mathcal{O}$, then we can calculate the following:

\vspace{20pt}
Fix $\epsilon = 1$. So if the SDG algorithm is 1-DP, then the following holds:
\begin{equation}
    \frac{P[\mathcal{A}(D) \in \mathcal{O}]}{P[\mathcal{A}(D') \in \mathcal{O}]} \leq e \approx 2.718
\end{equation}

Because this has to hold for any neighboring $D$ and $D'$, we can further say
\begin{equation}
    0.368 \approx \frac{1}{e} \leq \frac{P[\mathcal{A}(D) \in \mathcal{O}]}{P[\mathcal{A}(D') \in \mathcal{O}]} \leq e \approx 2.718
\end{equation}

I believe it is correct to say that $P[\mathcal{A}(D') \in \mathcal{O}]$ is the true positive rate (TPR), since we are saying $t \in D'$ and $\mathcal{A}$ yields a dataset in $\mathcal{O}$ that the attacker would correctly classify. Similarly, $P[\mathcal{A}(D) \in \mathcal{O}]$ is the false positive rate (FPR), since $t \notin D$, and so the attacker would incorrectly give a ``yes''. 

I'm going to try to calculate the theoretical maximum AUC an attacker can achieve. To do that, I will try to calculate the theoretical maximum TPR/FPR curves. First, if the attacker gives all ``yes''s, then the TPR is 100\%, and so the attacker's FPR must be
\begin{equation}
    0.368 \approx \frac{1}{e} \leq \frac{P[\mathcal{A}(D) \in \mathcal{O}] = FPR}{1} \leq e \approx 2.718
\end{equation}
so,
\begin{equation}
    FPR \geq \frac{1}{e} 
\end{equation}

By plugging in different TPRs, we get a linear lower bound for the FPR.

\begin{figure}[ht!]
\centering
\includegraphics[width=.46\textwidth]{./e1_minFPR.png}
\end{figure}


But what if we plug in different FPRs (for example, if the attacker always gives ``yes''s, then the FPR is 100\%)? Then we get the same boundary line. Take FPR = .2, for example.
\begin{equation}
    \frac{1}{e} \leq \frac{.2}{TPR} \leq e
\end{equation}
then,
\begin{equation}
    TPR \leq .2 e \approx .544
\end{equation}

\begin{figure}[ht!]
\centering
\includegraphics[width=.7\textwidth]{./e1_maxTPR.png}
\end{figure}

(I am leaving out the other side of the inequalities, i.e. when $FPR = .2$, $TPR \geq .2/e$. Or, when $FPR = 1$, $TPR \leq e$, because they yield boundary lines above 1.0 or below the ``random guess'' TPR = FPR line.)

Here is the TPR/FPR boundary for different $\epsilon$:

\begin{figure}[ht!]
\centering
\includegraphics[width=.7\textwidth]{./several_e.png}
\end{figure}

Surprisingly, for the reasonable setting when $\epsilon = 5$, the theoretical AUC is very very close to 1.0. But, as we expect, as $\epsilon$ approaches 0, the boundary approaches the random guessing line.


\newpage
Lastly, if these calculations are correct, then calculating maximum AUC is as simple as calculating the area of a triangle plus the area of a rectangle:

\begin{figure}[ht!]
\centering
\includegraphics[width=.7\textwidth]{./triangle_square.png}
\end{figure}


Here is what the maximum theoretical AUCs look like superimposed onto our MAMA-MIA attack results:


\begin{figure}[ht!]
\centering
\includegraphics[width=\textwidth]{./mama_mia_dp_boundary.png}
\caption{The yellow line is my``MAMA-MIA'' attack, and the other lines are the successes of other MIAs}
\end{figure}



\textbf{Note:} This analysis assumes pure DP, and does not yet calculate the bounds for $(\epsilon, \delta)$-DP or $(\alpha, \epsilon)$-DP. If my analysis above is correct, then it would be straightforward to calculate $(\epsilon, \delta)$-DP boundaries.




\end{document}